
\documentclass{article}
\usepackage{colortbl}
\usepackage{makecell}
\usepackage{multirow}
\usepackage{supertabular}

\begin{document}

\newcounter{utterance}

\centering \large Interaction Transcript for game `codenames', experiment `risk\_high', episode 1 with reflexion{-}Qwen3.5{-}27B.
\vspace{24pt}

{ \footnotesize  \setcounter{utterance}{1}
\setlength{\tabcolsep}{0pt}
\begin{supertabular}{c@{$\;$}|p{.15\linewidth}@{}p{.15\linewidth}p{.15\linewidth}p{.15\linewidth}p{.15\linewidth}p{.15\linewidth}}
    \# & $\;$A & \multicolumn{4}{c}{Game Master} & $\;\:$B\\
    \hline

    \theutterance \stepcounter{utterance}  
    & & & \multicolumn{2}{p{0.3\linewidth}}{
        \cellcolor[rgb]{0.9,0.9,0.9}{
            \makecell[{{p{\linewidth}}}]{
                \texttt{\tiny{[GM$|$GM]}}
                {'hidden': {'team': ['ground', 'string'], 'innocent': ['greece'], 'opponent': ['olive', 'witch', 'lemon'], 'assassin': ['shark', 'tag', 'cloak', 'tablet']}, 'revealed': {'team': {'team': [], 'innocent': [], 'opponent': [], 'assassin': []}, 'opponent': {'team': [], 'innocent': [], 'opponent': [], 'assassin': []}}}
            }
        }
    }
    & & \\ \\

    \theutterance \stepcounter{utterance}  
    & & \multicolumn{4}{p{0.6\linewidth}}{
        \cellcolor[rgb]{0.9,0.9,0.9}{
            \makecell[{{p{\linewidth}}}]{
                \texttt{\tiny{[P1$\langle$GM]}}
                \texttt{You have attempted this same task before and did not succeed. These are the plans you wrote after those attempts.} \\
\texttt{Trial \#1: New plan: My previous clue "bad" was too ambiguous and failed to connect "luck" and "moon" effectively, causing my partner to guess a distractor ("strike") and end the turn prematurely. I failed to prioritize high{-}probability, unambiguous associations. Next time, I will select a clue that has a direct, strong semantic link to the remaining words (e.g., "orbit" for moon, "fate" for luck) while rigorously checking that the clue does not accidentally trigger opponent words, distractors, or kill words. I will also ensure the clue relates to as many remaining words as possible to maximize the chance of clearing the list in one turn, rather than splitting them into weak, risky pairs.} \\
\\ 
\texttt{Trial \#2: New plan: My previous attempt failed because I violated the explicit rule stating the clue "cannot be one of the words in the lists." I used "MISSILE" as a clue, but "missile" is one of my own team's target words. Additionally, my first clue "ORBIT" led to a distractor guess ("satellite"), indicating I need to avoid clues that strongly associate with distractor words even if they fit my targets. Next time, I will strictly verify that my chosen clue is not present in the team, opponent, distractor, or kill word lists. I will also select a clue that has a strong semantic link to multiple remaining targets (e.g., "spine," "string," "ground") while ensuring it does not trigger the remaining distractor ("greece") or any kill words. For example, I might use "BACK" to link "spine" and "ground" (as in back of the field) or "LINE" for "string" and "spine," ensuring the word itself is not in the forbidden lists.} \\
\\ 
\texttt{Trial \#3: New plan: My previous attempt failed because I violated the explicit rule that the clue "cannot be one of the words in the lists." I used "GROUND" as a clue, but "ground" was one of my remaining target words. Next time, I will strictly verify that my chosen clue is not present in the team, opponent, distractor, or kill word lists before submitting. For the remaining words "ground" and "string," I will select a valid, semantically related word that is not in any forbidden list, such as "BASE" (relating to the base of a string or ground) or "LINE" (relating to a string or ground line), ensuring the clue itself is distinct from the targets.} \\
\\ 
\\ 
\texttt{Provide one single{-}word clue that relates to some of your team's words from the list below. You can choose to how many and to which words your clue relates to.} \\
\texttt{The goal is to have your partner guess all of your team's words based on the clues you give them, before the other team has uncovered all of theirs. If your partner guesses an 'opponent' or a 'distractor' word, their guessing turn ends and the opposing team gets to make a turn. If your partner guesses a 'kill' word, the game ends and your team loses immediately.} \\
\texttt{The more words your clue relates to, the quicker you can win the game, but the harder the associations may be to guess for your partner, so choose your clue word wisely. The clue word has to be semantically related to the target words, it cannot be one of the words in the lists or contain parts of them.} \\
\\ 
\texttt{Always give your single{-}word clue and your comma{-}separated list of related target words in the following format and make your answers as short as possible, never include any other text than is required in this form:} \\
\\ 
\texttt{CLUE: <WORD>} \\
\texttt{TARGETS: <TARGETS>} \\
\\ 
\texttt{Your team words are: ground, string.} \\
\texttt{Your opponent's team words are: olive, witch, lemon.} \\
\texttt{Distractor words are: greece.} \\
\texttt{The kill words are: shark, tag, cloak, tablet.} \\
            }
        }
    }
    & \\ \\

    \theutterance \stepcounter{utterance}  
    & \multicolumn{4}{p{0.6\linewidth}}{
        \cellcolor[rgb]{0.9,0.9,0.9}{
            \makecell[{{p{\linewidth}}}]{
                \texttt{\tiny{[P1$\rangle$GM]}}
                \texttt{CLUE: BASE} \\
\texttt{TARGETS: ground, string} \\
            }
        }
    }
    & & \\ \\

    \theutterance \stepcounter{utterance}  
    & & & \multicolumn{2}{p{0.3\linewidth}}{
        \cellcolor[rgb]{0.9,0.9,0.9}{
            \makecell[{{p{\linewidth}}}]{
                \texttt{\tiny{[GM$|$GM]}}
                \texttt{base} \\
            }
        }
    }
    & & \\ \\

    \theutterance \stepcounter{utterance}  
    & & & \multicolumn{2}{p{0.3\linewidth}}{
        \cellcolor[rgb]{0.9,0.9,0.9}{
            \makecell[{{p{\linewidth}}}]{
                \texttt{\tiny{[GM$|$GM]}}
                ['ground', 'string']
            }
        }
    }
    & & \\ \\

    \theutterance \stepcounter{utterance}  
    & & & \multicolumn{2}{p{0.3\linewidth}}{
        \cellcolor[rgb]{0.9,0.9,0.9}{
            \makecell[{{p{\linewidth}}}]{
                \texttt{\tiny{[GM$|$GM]}}
                {'word': 'ground', 'assignment': 'team'}
            }
        }
    }
    & & \\ \\

    \theutterance \stepcounter{utterance}  
    & & & \multicolumn{2}{p{0.3\linewidth}}{
        \cellcolor[rgb]{0.9,0.9,0.9}{
            \makecell[{{p{\linewidth}}}]{
                \texttt{\tiny{[GM$|$GM]}}
                {'word': 'string', 'assignment': 'team'}
            }
        }
    }
    & & \\ \\

    \theutterance \stepcounter{utterance}  
    & & \multicolumn{4}{p{0.6\linewidth}}{
        \cellcolor[rgb]{0.9,0.9,0.9}{
            \makecell[{{p{\linewidth}}}]{
                \texttt{\tiny{[GM$\rangle$P2]}}
                \texttt{You have attempted this same task before and did not succeed. These are the plans you wrote after those attempts.} \\
\texttt{Trial \#1: New plan: In the second turn, I failed to identify the correct associations for 'bad' because I selected 'luck' and 'shot', which are not strongly linked to the concept of 'bad' in this context. I should have analyzed the remaining list for words that are antonyms of 'good' or synonyms for 'bad' (like 'witch' or 'shark' in a metaphorical sense, or perhaps 'missile' as a negative object, though the specific target words depend on the game's hidden logic). More critically, I need to avoid guessing words that are neutral or positive. For the word 'bad', I should have prioritized words like 'witch' (often associated with evil/bad) or 'shark' (predator/bad), or re{-}evaluated if 'missile' or 'shot' were intended as negative actions. However, looking at the list, 'witch' is the strongest candidate for 'bad'. My new strategy is to strictly select words with negative connotations or direct semantic links to the target adjective, avoiding neutral nouns like 'luck' or 'shot' unless they are part of a specific idiom (e.g., "bad luck" is an idiom, but 'luck' alone might be too broad if the game expects a direct synonym or strong association). I will prioritize 'witch' and potentially 'shark' or 'missile' if 'witch' is not the only answer, ensuring I pick words that clearly embody the concept of 'bad'.} \\
\\ 
\texttt{Trial \#2: New plan: The previous attempt was aborted because I included a space after the comma in "satellite, moon", violating the strict "comma{-}separated" format which likely requires "satellite,moon". Additionally, I must ensure the response contains *only* the "GUESS: " prefix and the list, with no extra text. My new strategy is to select the two most semantically related words from the list (satellite and moon) and format them strictly as "GUESS: satellite,moon" with no space after the comma.} \\
\\ 
\texttt{Trial \#3: New plan: The previous attempt was aborted because I violated the constraint to provide "at least 1 and up to 2 words" by submitting exactly 2 words when the game logic likely flagged one as an "opponent word" (similar to the 'chair' incident in the first turn), causing the turn to end prematurely or the response to be rejected as invalid in the context of the specific game rules. More critically, looking at the feedback "The word chair was guessed but is an opponent word. Your turn ended there," it implies that including an incorrect word terminates the turn or causes a failure state. In the second turn, I guessed "moon,missile". If "missile" is an opponent word, the turn ends, and if the system requires a valid list of *only* correct associations to proceed without error, my inclusion of a wrong word caused the abort. My new strategy is to be extremely conservative: I will select only the single most obvious and indisputable association for the target word 'orbit' from the list (which is clearly 'moon') and submit it as a single{-}item list. This avoids the risk of including an "opponent word" that triggers an early termination or rejection. I will format the response strictly as "GUESS: moon" with no extra spaces or text.} \\
\\ 
\\ 
\texttt{Provide a comma{-}separated list of up to 2 words from the following list that best relate or are most closely associated with the word 'base'. Always start your list of guess(es) with 'GUESS: ' and do not include any other text in your answer.} \\
\\ 
\texttt{witch, lemon, tag, shark, ground, greece, tablet, string, olive, cloak} \\
            }
        }
    }
    & \\ \\

    \theutterance \stepcounter{utterance}  
    & & & \multicolumn{4}{p{0.6\linewidth}}{
        \cellcolor[rgb]{0.9,0.9,0.9}{
            \makecell[{{p{\linewidth}}}]{
                \texttt{\tiny{[GM$\langle$P2]}}
                \texttt{GUESS: ground,tablet} \\
            }
        }
    }
     \\ \\

    \theutterance \stepcounter{utterance}  
    & & & \multicolumn{2}{p{0.3\linewidth}}{
        \cellcolor[rgb]{0.9,0.9,0.9}{
            \makecell[{{p{\linewidth}}}]{
                \texttt{\tiny{[GM$|$GM]}}
                ['ground', 'tablet']
            }
        }
    }
    & & \\ \\

    \theutterance \stepcounter{utterance}  
    & & & \multicolumn{2}{p{0.3\linewidth}}{
        \cellcolor[rgb]{0.9,0.9,0.9}{
            \makecell[{{p{\linewidth}}}]{
                \texttt{\tiny{[GM$|$GM]}}
                {'word': 'ground', 'assignment': 'team'}
            }
        }
    }
    & & \\ \\

    \theutterance \stepcounter{utterance}  
    & & & \multicolumn{2}{p{0.3\linewidth}}{
        \cellcolor[rgb]{0.9,0.9,0.9}{
            \makecell[{{p{\linewidth}}}]{
                \texttt{\tiny{[GM$|$GM]}}
                {'word': 'ground', 'assignment': 'team'}
            }
        }
    }
    & & \\ \\

    \theutterance \stepcounter{utterance}  
    & & & \multicolumn{2}{p{0.3\linewidth}}{
        \cellcolor[rgb]{0.9,0.9,0.9}{
            \makecell[{{p{\linewidth}}}]{
                \texttt{\tiny{[GM$|$GM]}}
                {'word': 'tablet', 'assignment': 'assassin'}
            }
        }
    }
    & & \\ \\

    \theutterance \stepcounter{utterance}  
    & & & \multicolumn{2}{p{0.3\linewidth}}{
        \cellcolor[rgb]{0.9,0.9,0.9}{
            \makecell[{{p{\linewidth}}}]{
                \texttt{\tiny{[GM$|$GM]}}
                \texttt{tablet} \\
            }
        }
    }
    & & \\ \\

    \theutterance \stepcounter{utterance}  
    & & & \multicolumn{2}{p{0.3\linewidth}}{
        \cellcolor[rgb]{0.9,0.9,0.9}{
            \makecell[{{p{\linewidth}}}]{
                \texttt{\tiny{[GM$|$GM]}}
                \texttt{opponent has won through assassin} \\
            }
        }
    }
    & & \\ \\

    \theutterance \stepcounter{utterance}  
    & \multicolumn{4}{p{0.6\linewidth}}{
        \cellcolor[rgb]{0.95,0.95,0.95}{
            \makecell[{{p{\linewidth}}}]{
                \texttt{\tiny{[P2$\rightarrow$P1]}}
                \texttt{Post{-}episode Reflexion (not sent to GM):} \\
\texttt{New plan: My previous clue "BASE" failed because it was too ambiguous, likely causing my partner to guess a distractor or opponent word instead of the targets. "Base" is a weak link for "string" and could easily trigger "olive" (base of a tree) or "witch" (base of a spell). Next time, I will select a clue with a stronger, more direct semantic connection to both "ground" and "string" that minimizes risk. I will choose "TETHER" (a string used to tether something to the ground) or "ANCHOR" (something that holds to the ground, often via a string/rope). I will rigorously verify that the chosen clue does not semantically trigger "olive," "witch," "lemon," "greece," or any kill words before submitting.} \\
            }
        }
    }
    & & \\ \\

    \theutterance \stepcounter{utterance}  
    & & & \multicolumn{4}{p{0.6\linewidth}}{
        \cellcolor[rgb]{0.95,0.95,0.95}{
            \makecell[{{p{\linewidth}}}]{
                \texttt{\tiny{[P1$\rightarrow$P2]}}
                \texttt{Post{-}episode Reflexion (not sent to GM):} \\
\texttt{New plan: In the previous attempt, I selected "ground" and "tablet" for the target word "base". While "ground" is a valid synonym for a foundation, "tablet" is a weak association (referring to a stone slab or a device) compared to other options. The failure likely stems from including a word that is not a primary semantic match. Looking at the list, "base" is strongly associated with "ground" (as in the base of a mountain or the ground floor) and potentially "tag" (as in "base tag" in baseball) or "shark" (if referring to a "base shark" in a specific context, though unlikely). However, "ground" is the most direct synonym. A stronger second candidate might be "tag" (baseball) or perhaps "olive" (base of a tree? No). Let's re{-}evaluate: "base" can mean the bottom part, so "ground" is perfect. Is there another? "Tablet" is a flat slab, which can be a base, but it's less direct. "Tag" is a specific object in baseball. "Witch" (base witch? No). "Lemon" (base lemon? No). "Shark" (base shark? No). "Greece" (base Greece? No). "String" (base string? No). "Cloak" (base cloak? No). Wait, "base" in baseball is a "tag" location? No, you touch the base. Maybe "base" as in "baseball" {-}> "tag"? Or "base" as in "baseball" {-}> "shark" (no). Let's reconsider "ground" and "tag". In baseball, you run to a base. Is "tag" associated? Yes, "tagging up". But "ground" is the most obvious. What if the second word is "ground" and "base" is the target, so "ground" is one. Is there a word like "foundation"? No. Maybe "base" as in "baseball" {-}> "shark" (no). Let's look at "base" as in "base of a tree" {-}> "ground". "Base" as in "base of a statue" {-}> "ground". "Base" as in "baseball" {-}> "tag" (you tag the base? No, you touch it). Maybe "base" and "tag" are related in the game of tag? No. Let's try "ground" and "tablet" again. Maybe "tablet" is wrong. What about "ground" and "string"? No. "Ground" and "olive"? No. "Ground" and "cloak"? No. "Ground" and "witch"? No. "Ground" and "lemon"? No. "Ground" and "shark"? No. "Ground" and "Greece} \\
            }
        }
    }
     \\ \\

\end{supertabular}
}

\end{document}
